% =============================================================================
% Chapter 1 --- Introduction
% Grounded in README.md, Context 1 §1, and the 4 Background_Papers abstracts.
% =============================================================================

\chapter{Introduction}
\label{ch:introduction}

\section{Motivation: The Memory Wall and In-Memory Computing}
\label{sec:intro-motivation}

Modern computing is increasingly dominated by data-centric workloads, most
prominently artificial neural network (ANN) inference and training, in which
the energy and latency costs of moving data between memory and arithmetic
units dwarf the costs of the arithmetic itself. This imbalance is commonly
described as the \emph{memory wall}, and it arises in any architecture that
keeps storage and computation physically separate~\cite{Ielmini2018RRAM,
Verma2019IMC, Sebastian2020IMC}. For ANN workloads, whose core operation is
matrix--vector multiplication (MVM) over large weight matrices, the data
movement penalty becomes the dominant factor in achievable energy efficiency
and throughput.

\emph{In-memory computing} (IMC) is a family of architectural techniques
that addresses this bottleneck by performing certain computational tasks
directly inside a memory array, exploiting the physical attributes of the
memory devices together with their array-level
organization~\cite{Sebastian2020IMC}. For ANN acceleration, a 2D bit-cell
array is structurally aligned with the dataflow of MVM: weights are stored
in the array, inputs are applied on the rows (or columns), and the summed
analog current on the complementary lines constitutes the MVM
output~\cite{Verma2019IMC}. Non-volatile resistive switching devices such as
RRAM, PCM, and floating-gate memristors are natural candidates for the
storage elements of such arrays because their programmable conductance can
directly represent synaptic weights~\cite{Ielmini2018RRAM, Wang2022YFlash}.

\section{Project Context and Scope}
\label{sec:intro-context}

This report documents a course project in the Technion VLSI Laboratory
whose purpose is the design and functional verification of a
\emph{controller} for an ANN accelerator with in-memory computing
peripherals. The repository implements a controller-centric RTL stack in
which host traffic is decoded, sequenced into program, read, erase, and
inference operations, and driven to an ANN core via a packed 32-bit word
and a 3-bit pulse-mode bus. The controllable command tokens defined in the
parallel-interface package are \sv{CMD\_READ}, \sv{CMD\_PROG},
\sv{CMD\_ERASE}, \sv{CMD\_INF}, and \sv{CMD\_HIZ} (idle). The project scope
is explicitly confined to RTL implementation and functional simulation;
no netlist, timing closure, or power/area analysis artifacts are produced.

The design is motivated by a practical engineering observation that follows
directly from the physics of resistive-switching weight storage: a single
\emph{write once} operation is unreliable because of device-to-device and
cycle-to-cycle variability, so a closed-loop
\emph{program--verify--(re-program or erase)} control policy is needed to
drive stored weights toward their intended values~\cite{Ielmini2018RRAM,
Sebastian2020IMC, Wang2022YFlash}. The central piece of digital control in
such an accelerator is therefore a hierarchical finite-state machine that
mediates between the host and the analog core and that implements the
recovery logic.

\section{Contributions}
\label{sec:intro-contributions}

The concrete contributions recorded in this repository are:

\begin{itemize}
    \item Three synthesizable SystemVerilog modules: \modname{parallel\_interface}
          (combinational host decoder with one-hot tail validation),
          \modname{ann\_controller} (hierarchical FSM with PROG, VERIFY, and
          ERASE sub-FSMs and data-driven pulse shaping), and
          \modname{input\_buffer} (a 64-byte register-based store with
          byte-read and eight-lane bit-serial read paths).
    \item A shared package layer (\sv{controller\_pkg},
          \sv{parallel\_interface\_pkg}, \sv{input\_buffer\_pkg}) that
          centralizes address packing, one-hot predicates, pulse-train
          arithmetic, and LUT-based pulse shaping.
    \item A directed SystemVerilog verification environment with eleven
          testbenches, a Python-driven compile/simulate flow, wave wrappers
          and \sv{.do} scripts for GUI debugging, and human-readable report
          files written to a dedicated \sv{target/} directory.
    \item A behavioral memristor mock (shadow weight matrix + handshake
          generator) together with a fault-injection strategy that forces
          under- and over-read conditions during verify so that the
          recovery branches of the FSM are exercised in simulation.
\end{itemize}

\section{Report Organization}
\label{sec:intro-organization}

\Cref{ch:background} reviews the background of in-memory computing,
resistive-switching weight storage, and the write--verify problem that
motivates the controller's recovery FSM. \Cref{ch:architecture} gives a
system-level view of the three-module stack, the host-to-core word format,
and the command set. \Cref{ch:rtl} is the technical heart of the report:
each of the three modules is described in detail, including all main-FSM
states, the three sub-FSMs, the pulse-train arithmetic, and the LUT-based
programming path. \Cref{ch:verification} documents the verification
methodology: the testbench architecture, the shadow-matrix mock, the
fault-injection gating, and the per-testbench scoreboard structure.
\Cref{ch:results} reports the measured outcomes from the checked-in
simulation artifacts. \Cref{ch:analysis} interprets what the results prove
and \Cref{ch:limitations} catalogs the limitations of the present work
together with a concrete roadmap for follow-on verification work.
\Cref{ch:conclusion} closes the report. Seven appendices collect
diagrammatic, tabular, and source-level technical material that supports
the main chapters.
